\documentclass{amsart}
% For dealing with references we use the comment environment
\usepackage{verbatim}
\newenvironment{reference}{\comment}{\endcomment}
%\newenvironment{reference}{}{}
\newenvironment{slogan}{\comment}{\endcomment}
\newenvironment{history}{\comment}{\endcomment}

% For commutative diagrams we use Xy-pic
\usepackage[all]{xy}

% We use 2cell for 2-commutative diagrams.
\xyoption{2cell}
\UseAllTwocells

% We use multicol for the list of chapters between chapters
\usepackage{multicol}

% This is generally recommended for better output
\usepackage{lmodern}
\usepackage[T1]{fontenc}

% For cross-file-references

% Package for hypertext links:
\usepackage{hyperref}

% For any local file, say "hello.tex" you want to link to please

% Theorem environments.
%
\theoremstyle{plain}
\newtheorem{theorem}[subsection]{Theorem}
\newtheorem{proposition}[subsection]{Proposition}
\newtheorem{lemma}[subsection]{Lemma}

\theoremstyle{definition}
\newtheorem{definition}[subsection]{Definition}
\newtheorem{example}[subsection]{Example}
\newtheorem{exercise}[subsection]{Exercise}
\newtheorem{situation}[subsection]{Situation}

\theoremstyle{remark}
\newtheorem{remark}[subsection]{Remark}
\newtheorem{remarks}[subsection]{Remarks}

\numberwithin{equation}{subsection}

% Macros
%
\def\lim{\mathop{\mathrm{lim}}\nolimits}
\def\colim{\mathop{\mathrm{colim}}\nolimits}
\def\Spec{\mathop{\mathrm{Spec}}}
\def\Hom{\mathop{\mathrm{Hom}}\nolimits}
\def\Ext{\mathop{\mathrm{Ext}}\nolimits}
\def\SheafHom{\mathop{\mathcal{H}\!\mathit{om}}\nolimits}
\def\SheafExt{\mathop{\mathcal{E}\!\mathit{xt}}\nolimits}
\def\Sch{\mathit{Sch}}
\def\Mor{\mathop{\mathrm{Mor}}\nolimits}
\def\Ob{\mathop{\mathrm{Ob}}\nolimits}
\def\Sh{\mathop{\mathit{Sh}}\nolimits}
\def\NL{\mathop{N\!L}\nolimits}
\def\CH{\mathop{\mathrm{CH}}\nolimits}
\def\proetale{{pro\text{-}\acute{e}tale}}
\def\etale{{\acute{e}tale}}
\def\QCoh{\mathit{QCoh}}
\def\Ker{\mathop{\mathrm{Ker}}}
\def\Im{\mathop{\mathrm{Im}}}
\def\Coker{\mathop{\mathrm{Coker}}}
\def\Coim{\mathop{\mathrm{Coim}}}

% Boxtimes
%
\DeclareMathSymbol{\boxtimes}{\mathbin}{AMSa}{"02}

%
% Macros for moduli stacks/spaces
%
\def\QCohstack{\mathcal{QC}\!\mathit{oh}}
\def\Cohstack{\mathcal{C}\!\mathit{oh}}
\def\Spacesstack{\mathcal{S}\!\mathit{paces}}
\def\Quotfunctor{\mathrm{Quot}}
\def\Hilbfunctor{\mathrm{Hilb}}
\def\Curvesstack{\mathcal{C}\!\mathit{urves}}
\def\Polarizedstack{\mathcal{P}\!\mathit{olarized}}
\def\Complexesstack{\mathcal{C}\!\mathit{omplexes}}
% \Pic is the operator that assigns to X its picard group, usage \Pic(X)
% \Picardstack_{X/B} denotes the Picard stack of X over B
% \Picardfunctor_{X/B} denotes the Picard functor of X over B
\def\Pic{\mathop{\mathrm{Pic}}\nolimits}
\def\Picardstack{\mathcal{P}\!\mathit{ic}}
\def\Picardfunctor{\mathrm{Pic}}
\def\Deformationcategory{\mathcal{D}\!\mathit{ef}}

\setlength{\emergencystretch}{3em}
\begin{document}
\title[Stacks Project, Tag 0FEC]{455 --- Gysin classes compose for compatible exact sequences of virtual normal sheaves}
\maketitle
\noindent Copyright (C) 2005--2025 Johan de Jong. Stacks Project authors. Source: \href{https://stacks.math.columbia.edu/tag/0FEC}{Tag 0FEC}.

\noindent Editorial difficulty note (not author text). Difficulty H3: The proof must compare the normal-cone definitions for nested centres with excess normal directions. A simultaneous blowup reduces both centres to Cartier divisors, top-Chern factors remove the excess bundle, and a cartesian cone diagram identifies the two fundamental cone cycles under the remaining zero-section pullback..

\noindent Editorial provenance note (not author text). History: excerpt from revision \texttt{a0444\allowbreak 6e57e\allowbreak c1fbc\allowbreak 252a8\allowbreak 71afc\allowbreak ec775\allowbreak 2fb28\allowbreak 07b14}, chow.tex, lines 12324--12462. Reference presentation was adapted; original-author context and core support blocks were retained or added. The mathematical wording is unchanged. Licensed under GNU FDL 1.2 or later; no invariant sections or cover texts. See ../licenses/COPYING. Context blocks reproduce author definitions and supporting results; they are not additional counted problems.

\begin{definition}
\label{definition-bivariant-class}
\begin{reference}
Similar to \href{https://stacks.math.columbia.edu/bibliography}{Bibliography: F (Definition 17.1)}
\end{reference}
Let $(S, \delta)$ be as in Situation \ref{situation-setup}.
Let $f : X \to Y$ be a morphism of schemes locally of finite type over $S$.
Let $p \in \mathbf{Z}$.
A {\it bivariant class $c$ of degree $p$ for $f$} is given by a rule
which assigns to every locally of finite type morphism $Y' \to Y$
and every $k$ a map
$$
c \cap - : \CH_k(Y') \longrightarrow \CH_{k - p}(X')
$$
where $X' = Y' \times_Y X$, satisfying the following conditions
\begin{enumerate}
\item if $Y'' \to Y'$ is a proper, then
$c \cap (Y'' \to Y')_*\alpha'' = (X'' \to X')_*(c \cap \alpha'')$
for all $\alpha''$ on $Y''$ where $X'' = Y'' \times_Y X$,
\item if $Y'' \to Y'$ is flat locally of finite type of
fixed relative dimension, then
$c \cap (Y'' \to Y')^*\alpha' = (X'' \to X')^*(c \cap \alpha')$
for all $\alpha'$ on $Y'$, and
\item if $(\mathcal{L}', s', i' : D' \to Y')$ is as in
Definition \ref{definition-gysin-homomorphism}
with pullback $(\mathcal{N}', t', j' : E' \to X')$ to $X'$,
then we have $c \cap (i')^*\alpha' = (j')^*(c \cap \alpha')$
for all $\alpha'$ on $Y'$.
\end{enumerate}
The collection of all bivariant classes of degree $p$ for $f$ is
denoted $A^p(X \to Y)$.
\end{definition}

\section{Higher codimension gysin homomorphisms}
\label{section-gysin-higher-codimension}

\noindent
Let $(S, \delta)$ be as in Situation \ref{situation-setup}. Let $X$ be a scheme
locally of finite type over $S$. In this section we are going to consider
triples
$$
(Z \to X, \mathcal{N}, \sigma : \mathcal{N}^\vee \to \mathcal{C}_{Z/X})
$$
consisting of a closed immersion $Z \to X$ and a locally free
$\mathcal{O}_Z$-module $\mathcal{N}$ and a surjection
$\sigma : \mathcal{N}^\vee \to \mathcal{C}_{Z/X}$ from the dual
of $\mathcal{N}$ to the conormal sheaf of $Z$ in $X$, see
Morphisms, Section \href{https://stacks.math.columbia.edu/tag/01R1}{section conormal sheaf (Tag 01R1)}.
We will say
$\mathcal{N}$ is a {\it virtual normal sheaf for $Z$ in $X$}.



\begin{situation}
\label{situation-setup}
Here $S$ is a locally Noetherian, and universally catenary scheme.
Moreover, we assume $S$ is endowed with a dimension function
$\delta : S \longrightarrow \mathbf{Z}$.
\end{situation}

\begin{definition}
\label{definition-gysin-homomorphism}
Let $(S, \delta)$ be as in Situation \ref{situation-setup}.
Let $X$ be locally of finite type over $S$.
Let $(\mathcal{L}, s)$ be a pair consisting of an invertible
sheaf and a global section $s \in \Gamma(X, \mathcal{L})$.
Let $D = Z(s)$ be the zero scheme of $s$, and
denote $i : D \to X$ the closed immersion.
We define, for every integer $k$, a {\it Gysin homomorphism}
$$
i^* : Z_{k + 1}(X) \to \CH_k(D).
$$
by the following rules:
\begin{enumerate}
\item Given an integral closed subscheme $W \subset X$ with
$\dim_\delta(W) = k + 1$ we define
\begin{enumerate}
\item if $W \not \subset D$, then $i^*[W] = [D \cap W]_k$ as a
$k$-cycle on $D$, and
\item if $W \subset D$, then
$i^*[W] = i'_*(c_1(\mathcal{L}|_W) \cap [W])$,
where $i' : W \to D$ is the induced closed immersion.
\end{enumerate}
\item For a general $(k + 1)$-cycle $\alpha = \sum n_j[W_j]$
we set
$$
i^*\alpha = \sum n_j i^*[W_j]
$$
\item If $D$ is an effective Cartier divisor, then we denote
$D \cdot \alpha = i_*i^*\alpha$ the pushforward of the class $i^*\alpha$
to a class on $X$.
\end{enumerate}
\end{definition}

\noindent
Let $(S, \delta)$ be as in Situation \ref{situation-setup}. Let $X$ be a scheme
locally of finite type over $S$. Let $\mathcal{N}$ be a virtual normal bundle
for a closed immersion $Z \to X$. In this situation we set
$$
p : N = \underline{\Spec}_Z(\text{Sym}(\mathcal{N}^\vee)) \longrightarrow Z
$$
equal to the vector bundle over $Z$
whose sections correspond to sections of $\mathcal{N}$.
In this situation we have canonical closed immersions
$$
C_ZX \longrightarrow N_ZX \longrightarrow N
$$
The first closed immersion is Divisors, Equation
(\href{https://stacks.math.columbia.edu/tag/0637}{equation normal cone in normal bundle (Tag 0637)})
and the second closed immersion corresponds to the surjection
$\text{Sym}(\mathcal{N}^\vee) \to \text{Sym}(\mathcal{C}_{Z/X})$
induced by $\sigma$.
Let
$$
b : W \longrightarrow \mathbf{P}^1_X
$$
be the blowing up in $\infty(Z)$ constructed in
Section \href{https://stacks.math.columbia.edu/tag/0FBG}{section blowup Z first (Tag 0FBG)}. By
Lemma \href{https://stacks.math.columbia.edu/tag/0F9J}{lemma gysin at infty (Tag 0F9J)}
we have a canonical bivariant class in
$$
C \in A^0(W_\infty \to X)
$$
Consider the open immersion $j : C_ZX \to W_\infty$ of
(\href{https://stacks.math.columbia.edu/tag/0FBH}{item cone is open (Tag 0FBH)}) and the closed immersion
$i : C_ZX \to N$ constructed above. By Lemma \href{https://stacks.math.columbia.edu/tag/02TY}{lemma vectorbundle (Tag 02TY)}
for every $\alpha \in \CH_k(X)$ there exists a unique
$\beta \in \CH_*(Z)$ such that
$$
i_*j^*(C \cap \alpha) = p^*\beta
$$
We set $c(Z \to X, \mathcal{N}) \cap \alpha = \beta$.


\begin{lemma}
\label{lemma-gysin-composition}
Let $(S, \delta)$ be as in Situation \ref{situation-setup}.
Let $Z \subset Y \subset X$ be closed subschemes of a scheme locally
of finite type over $S$.
Let $\mathcal{N}$ be a virtual normal sheaf for $Z \subset X$.
Let $\mathcal{N}'$ be a virtual normal sheaf for $Z \subset Y$.
Let $\mathcal{N}''$ be a virtual normal sheaf for $Y \subset X$.
Assume there is a commutative diagram
$$
\xymatrix{
(\mathcal{N}'')^\vee|_Z \ar[r] \ar[d] &
\mathcal{N}^\vee \ar[r] \ar[d] &
(\mathcal{N}')^\vee \ar[d] \\
\mathcal{C}_{Y/X}|_Z \ar[r] &
\mathcal{C}_{Z/X} \ar[r] &
\mathcal{C}_{Z/Y}
}
$$
where the sequence at the bottom is from More on Morphisms, Lemma
\href{https://stacks.math.columbia.edu/tag/06AE}{lemma transitivity conormal (Tag 06AE)} and the top
sequence is a short exact sequence. Then
$$
c(Z \to X, \mathcal{N}) =
c(Z \to Y, \mathcal{N}') \circ c(Y \to X, \mathcal{N}'')
$$
in $A^*(Z \to X)^\wedge$.
\end{lemma}

\begin{proof}
Observe that the assumptions remain satisfied after any base change
by a morphism $X' \to X$ which is locally of finite type (the short
exact sequence of virtual normal sheaves is locally split hence
remains exact after any base change). Thus to check the
equality of bivariant classes we may use Lemma \href{https://stacks.math.columbia.edu/tag/02UC}{lemma bivariant zero (Tag 02UC)}.
Thus we may assume $X$ is an integral scheme and we have to show
$c(Z \to X, \mathcal{N}) \cap [X] =
c(Z \to Y, \mathcal{N}') \cap c(Y \to X, \mathcal{N}'') \cap [X]$.

\medskip\noindent
If $Y = X$, then we have
\begin{align*}
c(Z \to Y, \mathcal{N}') \cap c(Y \to X, \mathcal{N}'') \cap [X]
& =
c(Z \to Y, \mathcal{N}') \cap c_{top}(\mathcal{N}'') \cap [Y] \\
& =
c_{top}(\mathcal{N}''|_Z) \cap c(Z \to Y, \mathcal{N}') \cap [Y] \\
& =
c(Z \to X, \mathcal{N}) \cap [X] 
\end{align*}
The first equality by Lemma \href{https://stacks.math.columbia.edu/tag/0FBL}{lemma gysin decompose (Tag 0FBL)}.
The second because Chern classes commute with bivariant classes
(Lemma \href{https://stacks.math.columbia.edu/tag/02UA}{lemma cap commutative chern (Tag 02UA)}).
The third equality by Lemma \href{https://stacks.math.columbia.edu/tag/0FBL}{lemma gysin decompose (Tag 0FBL)}.

\medskip\noindent
Assume $Y \not = X$. By Lemma \href{https://stacks.math.columbia.edu/tag/02UC}{lemma bivariant zero (Tag 02UC)}
it even suffices to prove the result after blowing up $X$ in a nonzero ideal.
Let us blowup $X$ in the product of the ideal sheaf of $Y$ and the ideal
sheaf of $Z$. This reduces us to the case where both $Y$ and $Z$ are
effective Cartier divisors on $X$, see
Divisors, Lemmas
\href{https://stacks.math.columbia.edu/tag/02OS}{lemma blowing up gives effective Cartier divisor (Tag 02OS)} and
\href{https://stacks.math.columbia.edu/tag/080A}{lemma blowing up two ideals (Tag 080A)}.

\medskip\noindent
Denote $\mathcal{N}'' \to \mathcal{E}$ the surjection of finite locally
free $\mathcal{O}_Z$-modules such that
$0 \to \mathcal{E}^\vee \to (\mathcal{N}'')^\vee \to \mathcal{C}_{Y/X} \to 0$
is a short exact sequence. Then $\mathcal{N} \to \mathcal{E}|_Z$
is a surjection as well. Denote $\mathcal{N}_1$ the finite locally free kernel
of this map and observe that $\mathcal{N}^\vee \to \mathcal{C}_{Z/X}$
factors through $\mathcal{N}_1$.
By Lemma \href{https://stacks.math.columbia.edu/tag/0FBL}{lemma gysin decompose (Tag 0FBL)} we have
$$
c(Y \to X, \mathcal{N}'') = c_{top}(\mathcal{E}) \circ
c(Y \to X, \mathcal{C}_{Y/X}^\vee)
$$
and
$$
c(Z \to X, \mathcal{N}) = c_{top}(\mathcal{E}|_Z) \circ
c(Z \to X, \mathcal{N}_1)
$$
Since Chern classes of bundles commute with bivariant classes
(Lemma \href{https://stacks.math.columbia.edu/tag/02UA}{lemma cap commutative chern (Tag 02UA)})
it suffices to prove
$$
c(Z \to X, \mathcal{N}_1) =
c(Z \to Y, \mathcal{N}') \circ c(Y \to X, \mathcal{C}_{Y/X}^\vee)
$$
in $A^*(Z \to X)$. This we may assume that $\mathcal{N}'' = \mathcal{C}_{Y/X}$.
This reduces us to the case discussed in the next paragraph.

\medskip\noindent
In this paragraph $Z$ and $Y$ are effective Cartier divisors on $X$
integral of dimension $n$, we have $\mathcal{N}'' = \mathcal{C}_{Y/X}$.
In this case $c(Y \to X, \mathcal{C}_{Y/X}^\vee) \cap [X] = [Y]_{n - 1}$ by
Lemma \href{https://stacks.math.columbia.edu/tag/0FBM}{lemma gysin fundamental (Tag 0FBM)}. Thus we have to prove that
$c(Z \to X, \mathcal{N}) \cap [X] = c(Z \to Y, \mathcal{N}') \cap [Y]_{n - 1}$.
Denote $N$ and $N'$ the vector bundles over $Z$ associated to
$\mathcal{N}$ and $\mathcal{N}'$. Consider the commutative diagram
$$
\xymatrix{
N' \ar[r]_i &
N \ar[r] &
(C_Y X) \times_Y Z \\
C_Z Y \ar[r] \ar[u] &
C_Z X \ar[u]
}
$$
of cones and vector bundles over $Z$. Observe that $N'$ is a relative
effective Cartier divisor in $N$ over $Z$ and that
$$
\xymatrix{
N' \ar[d] \ar[r]_i & N \ar[d] \\
Z \ar[r]^-o & (C_Y X) \times_Y Z
}
$$
is cartesian where $o$ is the zero section of the line bundle
$C_Y X$ over $Y$. By
Lemma \href{https://stacks.math.columbia.edu/tag/0FEB}{lemma relation normal cones (Tag 0FEB)} we have $o^*[C_ZX]_n = [C_Z Y]_{n - 1}$
in
$$
\CH_{n - 1}(Y \times_{o, C_Y X} C_ZX) =
\CH_{n - 1}(Z \times_{o, (C_Y X) \times_Y Z} C_ZX)
$$
By the cartesian property of
the square above this implies that
$$
i^*[C_ZX]_n = [C_Z Y]_{n - 1}
$$
in $\CH_{n - 1}(N')$. Now observe that
$\gamma = c(Z \to X, \mathcal{N}) \cap [X]$ and
$\gamma' = c(Z \to Y, \mathcal{N}') \cap [Y]_{n - 1}$
are characterized by $p^*\gamma = [C_Z X]_n$ in $\CH_n(N)$
and by $(p')^*\gamma' = [C_Z Y]_{n - 1}$ in $\CH_{n - 1}(N')$.
Hence the proof is finished as $i^* \circ p^* = (p')^*$ by
Lemma \href{https://stacks.math.columbia.edu/tag/02TR}{lemma relative effective cartier (Tag 02TR)}.
\end{proof}
\end{document}
